\documentclass[10pt,a4paper]{article}
\usepackage{amsmath,amssymb,amsthm,graphicx}
% Card-local minimal preamble for S7 cards cut from sources/markov-quantile-jep-2022-177/source0.tex
% (doi:10.5802/jep.177, JEP 9 (2022), CC BY 4.0). Author text untouched.
% Omitted and why: the publisher cedram class (unavailable; assembler supplies
% article) and the publisher date commands (\datereceived/\dateaccepted/
% \dateepreuves/\datepublished), per the unavailable-class precedent.
\usepackage{amsmath,amssymb,amsthm}
\usepackage[scr=rsfs,cal=euler]{mathalfa}
\usepackage{xfrac}
\multlinegap0pt
\renewcommand{\theenumi}{\alph{enumi}}
\renewcommand{\theenumii}{\roman{enumii}}
\newenvironment{enumeratei}
{\bgroup\def\theenumi{\roman{enumi}}\def\theenumii{\arabic{enumii}}\begin{enumerate}}
{\end{enumerate}\egroup}
\makeatletter
\let\p@enumii\relax
\makeatother
\newcommand{\ldot}{\mathbin{.}}
\newcommand{\lcdots}{\ldot\cdots\ldot}
\newcommand{\RedefinitSymbole}[1]{%
\expandafter\let\csname old\string#1\endcsname=#1
\let#1=\relax
\newcommand{#1}{\csname old\string#1\endcsname\,}%
}
\RedefinitSymbole{\forall} \RedefinitSymbole{\exists}
\newcommand\mto{\mathchoice{\longmapsto}{\mapsto}{\mapsto}{\mapsto}}
\newcommand{\psfrac}[2]{\sfrac{(#1)}{#2}}
\newcommand{\spfrac}[2]{\sfrac{#1}{(#2)}}
\usepackage[all,cmtip]{xy}\let\labelstyle\textstyle
\usepackage{stmaryrd}
\usepackage{bbm}
\usepackage{mathtools}
\theoremstyle{plain}
\newtheorem{themBJ}{Theorem}
\renewcommand*{\thethemBJ}{\Alph{themBJ}}
\newtheorem{them}{Theorem}[section]
\newtheorem{pro}[them]{Proposition}
\newtheorem{pronotation}[them]{Proposition/Notation}
\newtheorem{prodefinition}[them]{Proposition/Definition}
\newtheorem{lemnotation}[them]{Lemma/Notation}
\newtheorem{cor}[them]{Corollary}
\newtheorem{lem}[them]{Lemma}
\newtheorem*{claim}{Claim}
\theoremstyle{definition}
\newtheorem{vocabulary}[them]{Vocabulary}
\newtheorem{defi}[them]{Definition}
\newtheorem{convention}[them]{Convention}
\newtheorem{notation}[them]{Notation}
\newtheorem{remnotation}[them]{Remark/Notation}
\newtheorem{definotation}[them]{Definition/Notation}
\newtheorem{remdefi}[them]{Remark/Definition}
\newtheorem{reminder-notation}[them]{Reminder/Notation}
\newtheorem{rem}[them]{Remark}
\newtheorem{ex}[them]{Example}
\newtheorem{reminder}[them]{Reminder}
\newtheorem*{problem*}{Main Problem}
\newtheorem*{note}{Note}
\newtheorem*{fact}{Fact}
\newcommand{\dd}{\mathrm{d}}
\newcommand{\rec}{\mathfrak R}
\newcommand{\id}{\mathrm{Id}}
\newcommand{\join}{\mathrm{Joint}}
\newcommand{\como}{\mathfrak{Q}}
\newcommand{\mq}{\mathfrak{MQ}}
\newcommand{\lev}{\mathrm{Lev}}
\newcommand{\ML}{\mathrm{LK}}
\newcommand{\IT}{\mathrm{IK}}
\newcommand{\disp}{\mathfrak{Disp}}
\newcommand{\argth}{\operatorname{Argth}}
\newcommand{\vol}{\operatorname{Vol}}
\newcommand{\parti}{\operatorname{Part}}
\newcommand{\idk}{\operatorname{id}}
\newcommand{\espace}{E}
\newcommand{\subd}{R}
\newcommand{\Tt}{{\mathcal T}}
\newcommand{\XX}{{\mathcal X}}
\newcommand{\EE}{{\mathcal E}}
\newcommand{\A}{{\mathcal A}}
\newcommand{\Var}{{\mathcal L}}
\newcommand{\Id}{\operatorname{Id}}
\newcommand{\one}{\mathbbm{1}}
\let\mb\relax
\let\mn\relax
\newcommand{\trans}{{}^{\mathrm{t}}\mkern-1mu}
\newcommand{\eqdef}{\coloneqq}
\newcommand{\eqdefup}{\coloneqq}
\newcommand{\espcdot}{\mbox{$\,\cdot\,$}}
\let\eps\varepsilon
\newcommand{\Lg}{\lambda|_{[0,1]}}
\newcommand{\la}{\lambda}
\renewcommand{\P}{\mathbb{P}}
\newcommand{\p}{\mathcal{P}}
\newcommand{\m}{\mathcal{M}}
\newcommand{\ma}{\mathrm{Marg}}
\newcommand{\cC}{\mathcal{C}}\let\C\cC
\newcommand{\R}{\mathbb{R}}
\newcommand{\Q}{\mathbb{Q}}
\newcommand{\N}{\mathbb{N}}
\newcommand{\Z}{\mathbb{Z}}
\newcommand{\AND}{\quad\text{and}\quad}
\newcommand{\E}{\mathbb{E}}
\newcommand{\law}{\operatorname{Law}}
\newcommand{\var}{\operatorname{Var}}
\newcommand{\sto}{\mathrm{sto}}
\newcommand{\stosup}{\operatorname{stosup}}
\newcommand{\stoinf}{\operatorname{stoinf}}
\newcommand{\lcsup}{\operatorname{losup}}
\newcommand{\lcinf}{\operatorname{loinf}}
\newcommand{\leqp}{\preceq_+}
\newcommand{\leqc}{\preceq_{C}}
\newcommand{\leqcs}{\preceq_{C,\sto}}
\newcommand{\leqs}{\preceq_{\text{\rm sto}}}
\newcommand{\geqs}{\succeq_{\text{\rm sto}}}
\newcommand{\e}{\mathrm{e}}
\newcommand{\pr}{\operatorname{proj}}
\newcommand{\pdec}{\mathcal{P}^{\scriptscriptstyle\searrow}}
\newcommand{\mdec}{\mathcal{M}^{\scriptscriptstyle\searrow}}
\newcommand{\leb}{\LL_{[0,1]}}
\newcommand{\op}{\mathopen}
\newcommand{\clo}{\mathclose}
\newcommand{\leqlc}{\preceq_{\rm lo}}
\newcommand{\geqlc}{\succeq_{\rm lo}}
\newcommand{\rsemibracket}{\mathclose{\rfloor\hspace*{-.9ex}\lceil}}
\newcommand{\lsemibracket}{\mathclose{\rfloor\hspace*{-.9ex}\lceil}}
\newcommand{\semi}{\mathclose{\rfloor\hspace*{-.8ex}\lceil}}
\newlength{\arraycolsepsauvegardegenerale}
\setlength{\arraycolsepsauvegardegenerale}{\arraycolsep}
\newlength{\arraycolsepsauvegardetemporaire}
\newlength{\tabcolsepsauvegardegenerale}
\setlength{\tabcolsepsauvegardegenerale}{\tabcolsep}
\newlength{\tabcolsepsauvegardetemporaire}
\newenvironment{nnarray}%
{\setlength{\arraycolsepsauvegardetemporaire}{\arraycolsep}\setlength{\arraycolsep}{0cm}\begin{array}}
{\end{array}\setlength{\arraycolsep}{\arraycolsepsauvegardetemporaire}}
\newenvironment{nntabular}%
{\setlength{\tabcolsepsauvegardetemporaire}{\tabcolsep}\setlength{\tabcolsep}{0cm}\begin{tabular}}
{\end{tabular}\setlength{\tabcolsep}{\tabcolsepsauvegardetemporaire}}
\providecommand{\LL}{\mathbf{L}}
\providecommand{\ie}{i.e.}
\usepackage{cleveref}
\crefname{originaltheorem}{original theorem}{original theorems}
\Crefname{originaltheorem}{Original theorem}{Original theorems}
\makeatletter
\newlabel{defi:martingale}{{1.15}{source}{}{sourcepointer-1}{}}
\newlabel{defi:martingale@cref}{{[originaltheorem][0][]1.15}{[1][0][]source}{}{}{}}
\makeatother
\crefname{originaltheorem}{original theorem}{original theorems}
\Crefname{originaltheorem}{Original theorem}{Original theorems}
\makeatletter
\newlabel{item:increasing_coupling}{{1.15}{source}{}{sourcepointer-2}{}}
\newlabel{item:increasing_coupling@cref}{{[originaltheorem][0][]1.15}{[1][0][]source}{}{}{}}
\makeatother
\crefname{originaltheorem}{original theorem}{original theorems}
\Crefname{originaltheorem}{Original theorem}{Original theorems}
\makeatletter
\newlabel{rem:stosup_existe}{{3.9}{source}{}{sourcepointer-3}{}}
\newlabel{rem:stosup_existe@cref}{{[originaltheorem][0][]3.9}{[1][0][]source}{}{}{}}
\makeatother
\title{If P is a probability measure on with increasing couplings (see Definition ), there exists...}
\author{Boubel, Charles, Juillet, Nicolas}
\date{doi:10.5802/jep.177; CC BY 4.0}
\begin{document}
\maketitle
\setcounter{section}{1}
\setcounter{them}{18}
\setcounter{equation}{4}
\setcounter{themBJ}{3}
\noindent\textit{Editorial note.} All delivered mathematical text is the author's own, copied from the frozen source sources/markov-quantile-jep-2022-177/source0.tex; no statement or proof step was rewritten, reordered or removed. References inside the span are handed to the assembler as navigation stubs; load-bearing status is recorded per stub.\par\medskip
\noindent\textit{Editorial scope.} Single lem card cut from doi:10.5802/jep.177 (The Markov-quantile process attached to a family of marginals); statement and proof are the author own text at source lines 460-461 and 1301-1321.\par\medskip
\ifdefined\hypertarget\hypertarget{sourcepointer-1}{}\fi
\noindent\textit{Source reference.} Original theorem 1.15: sources/markov-quantile-jep-2022-177/source0.tex, source lines 405--405. navigation stub only: the target lives outside the counted span in the same article; whether it is load-bearing is judged per occurrence\par\medskip
\ifdefined\hypertarget\hypertarget{sourcepointer-2}{}\fi
\noindent\textit{Source reference.} Original theorem 1.15: sources/markov-quantile-jep-2022-177/source0.tex, source lines 415--415. navigation stub only: the target lives outside the counted span in the same article; whether it is load-bearing is judged per occurrence\par\medskip
\ifdefined\hypertarget\hypertarget{sourcepointer-3}{}\fi
\noindent\textit{Source reference.} Original theorem 3.9: sources/markov-quantile-jep-2022-177/source0.tex, source lines 1034--1034. navigation stub only: the target lives outside the counted span in the same article; whether it is load-bearing is judged per occurrence\par\medskip
\begin{lem}\label{coro}If $P$ is a probability measure on $\R^\R$ with increasing couplings (see Definition \ref{defi:martingale}\eqref{item:increasing_coupling}), there exists an increasing process $X=(X_t)_{t\in\R}:\Omega\to\R$, \ie the functions $t\mto X_t(\omega)$ are increasing, such that $P=\law(X)$.
\end{lem}

\begin{proof}[Proof of Lemma \ref{coro}]
Let $(\mu_t)_{t\in \R}$ be an increasing family of probability measures for $\leqs$ and $P$ be a Markov measure in $\ma((\mu_t)_{t\in \R})$ such that for every $S=\{s_1,\ldots,s_d\}$, the measure $(\pr^S)_\#P$ is concentrated on
\[
\{(x_1,\ldots,x_d)\in \R^d:\,x_1\leq\ldots \leq x_d\}.
\]
Set $\mu_{t^-}\eqdef\stosup_{s<t}\mu_s$ and $\mu_{t^+}\eqdef\stoinf_{s>t}\mu_s$, that are also the left and right limits of $(\mu_t)_t$ for the weak topology, see Remark \ref{rem:stosup_existe}. Since $(\mu_t)_{t\in \R}$ is increasing for $\leqs$, we have $\mu_{t^-}\leqs \mu_t \leqs \mu_{t^+}$ and $t\in\R$ is a discontinuity time of $(\mu_t)_{t\in \R}$ for the weak topology if and only if $\mu_{t^-}\neq\mu_{t^+}$. Such points are at most countably many. Indeed, the regions in $\R^2$ between the graphs of $F[\mu_{t^-}]$ and $F[\mu_{t^+}]$, for the discontinuity times~$t$, are disjoint and of positive Lebesgue measure, hence are at most countably many, by $\sigma$-additivity of the measure.

Let $C$ be a countable dense subset of $\R$ containing the discontinuity points. Introduce:\vspace*{-3pt}
\[N=\{x\in \R^\R:\,\exists (s,t)\in C^2, s<t\text{ and }x(s)>x(t)\}
.\]
Being a countable union of $P$-null sets, $N$ is $P$-null. Now take $(X_t)_{t\in \R}$ with law $P$, e.g., take the canonical process $\Omega\eqdef\R^\R$ and $X_t=\pr^t:x\in \R^\R\mto x(t)$. We define $(\tilde X_t)_{t}$ as null functions on $N$, and as follows on $\Omega\setminus N$:
\begin{itemize}
\item
for $t\in C$, $\tilde X_t(\omega)=X_t(\omega)$,

\item
for $t\notin C$, $\tilde X_t(\omega)=\lim_{s<t,\,s\in C}X_s(\omega)$.

\end{itemize}
Hence for every $\omega\in \Omega$, the curve $t\in C\mto X_t(\omega)$ is increasing and, even better, $t\in \R\mto X_t(\omega)$ is increasing. We are left to prove that $X_t=\tilde X_t$ almost surely. This is clear for $t\in C$. For each $t\in\R\setminus C$, $\{\omega\in\Omega:\exists s\in C, s<t\ \text{and}\ X_s(\omega)>X_t(\omega)\}$ is a union of null sets, hence is null, so almost surely, $X_t\geqslant \sup_{s<t,\,s\in C}X_s=\tilde X_t$. Besides $X_s\to_{s<t,\,s\in C, s\to t}\tilde{X}_t$ almost surely and thus in law, so that $\law(\tilde X_t)=\mu_{t^-}$. Moreover, $t$~is a continuity point of $(\mu_t)_t$, so that $\law(\tilde X_t)=\mu_{t}=\law(X_t)$. Thus, $X_t=\tilde{X}_t$ almost surely.
\end{proof}

\end{document}
